% !TeX root = ../papers/nc-p3.tex
\section{三薄色の捻りの無印同型と他枝との交差}
\label{sec:three-lines-twisted-isomorphism}
\begin{theorem}
定理\ref{thm:klein-three-lines}の捻った表示では、四係数 $p_i$ の多重集合の一致が添字保持同型の必要十分条件である。
二対反復の場合にも追加同一視はない。三薄色の二族、定理\ref{thm:transverse-doublet-isomorphism}の異方向族、
および一方の切断だけに欠色のある正則曲線は、無印の同型類として互いに交わらない。
\end{theorem}
\begin{proof}[関係空間の低階数曲線]
捻りを二対の内部と対の間の符号に分け、第 $i$ 行に $(-1)^{i_{\rm hi}i_{\rm lo}}$ を掛ける。
隣接関係空間 $R$ は次の行列の空間になる。
\[
 \begin{pmatrix}A&x&y&z\\x&B&z&y\\y&-z&C&-x\\-z&y&-x&D\end{pmatrix},
 \qquad p_0 A+p_1 B+p_2 C+p_3 D=0 .
\]
以下その対角を $d_i$、混合成分を $\xi_1=x,\xi_2=y,\xi_3=z$、対角部分空間を $D_p$ と書く。
階数2以下の射影閉集合は左右独立基底変更の不変量である。
$i,j$ の補対を $k,l$、$h=i+j\in K$ とすると、随伴行列の $(i,j),(j,i)$ 成分の符号付き和差の半分は、
各々符号を除いて
\[
 \xi_h(x^2+y^2-z^2-d_kd_l),\qquad
 \Big(\prod_{b\ne h,\, b>0}\xi_b\Big)(d_k-d_l)
\]
である。全 $\xi_b\ne0$ なら対角は定数で、$\sum p_i=1$ より零となり、残る点は平面二次曲線上にある。
二成分だけが非零なら、消えた色の二対で対角がそれぞれ一致する。二つの対角値には非自明な一次条件があり、
射影平面内でさらに上の非零な二次式が課される。従っていずれも正種数の成分を持たない。
三成分とも零ならちょうど六つの階数2点である。

一成分 $\xi_h$ だけが非零なら行列は二つの二行ブロックに分かれ、双方の行列式が零になる。
$h=3$ では $\xi_h$ を $\sqrt{-1}$ 倍して、$d_i d_j=d_kd_l=\xi_h^2$ と書ける。
$\xi_h=1$ のアフィン部分は $d_i=s,d_j=1/s,d_k=t,d_l=1/t$ により
\[
 \lambda=p_i s+p_j/s=-(p_k t+p_l/t)
\]
となる。二つの二次被覆の分岐対は $\{\pm 2\sqrt{p_i p_j}\},\{\pm 2\sqrt{p_k p_l}\}$ である。
対が異なれば交わらず、合成体の次数は4、分岐指数2の点が計八つで、正規化の種数は1となる。
この曲線の張る空間は混合成分を一つ許した $\mathbb P^3=P_h$ 全体である。実際 $1,\lambda$ と
二つの相異なる判別式の平方根が定数体上一次独立で、$1,s,s^{-1},t,t^{-1}$ の間の一次関係は表示した一つだけである。
分岐対が等しければ二つの二次体は一致し、成分は有理的である。
以上で階数2以下の集合の次元は高々1、正種数成分の個数は相補対積の異なる分割数と分かった。
\end{proof}

\begin{proof}[楕円曲線が複数ある場合]
二つの相補対積の等式、例えば $p_0p_1=p_2p_3$, $p_0p_2=p_1p_3$ から
$p_1=\pm p_2$, $p_0=\pm p_3$ で符号は共通となる。負号は成分和1に反する。
従って楕円曲線が一本なら二値の二対、零本なら全係数が等しく一つの類だけである。
二本以上ならその張る空間の交わりは $\mathbb P D_p$ であり、色を指定せず復元できる。
内部の六点の像と行空間は両側の四文字から二つを選ぶ六平面である。相異なるそれらの一次元の交わりから
四つの文字直線が決まる。像と行空間で同じ対が組になるため、隣接位置の文字置換は一致する。
四位置に適用すると全位置共通の置換と位置別対角変更だけが可能である。

全ての四文字置換は四元群のアフィン変換である。平行移動の捻り符号への寄与は位置別指標と共通定数に分かれる。
線形変換では交代化が二元体上の唯一の非零交代形式なので、元との比は対称双指標になる。
これは代数閉体上でコバウンダリである（基底の二文字で対角符号の平方根を選び、和での値を乗法則で定める）。
総色零の四語では従って各位置の倍率で吸収できる。よって全ての置換が実現する。

置換を除いた倍率を $f_j(i)$ とする。非純語は全て同じ共通倍率を受ける必要がある。
$(i,i,k,k)$ の全順列を比較すると $f_j(i)/f_j(k)$ は位置に依らない。そのため位置別定数を除き倍率は一つの関数 $f(i)$ である。
全 $i\ne k$ で $f(i)^2 f(k)^2$ が一定だから四平方は等しく、純語も非純語と同じ共通倍率になる。
従って $1-1/p_i$ が一致し、これには純成分が零の値も含まれる。
\end{proof}

\begin{proof}[一本の楕円曲線]
$p=(a,a,b,b)$、$r=(a-b)/(a+b)\ne0,\pm1$ とする。一意な楕円曲線の張る $P_1$ では $y=z=0$ である。
その中の階数2以下の点は楕円曲線の閉包と二つの孤立点からなる。
実際 $x=0$ の六点のうち、二つのブロックから一文字ずつ取る四点が二行列式の共通零点で、上の曲線の無限遠点である
（各被覆の $\lambda=\infty$ の二点は非分岐）。他の二点だけが一つのブロックに支えられる。
孤立点の像と行空間が左右それぞれに二平面の分解を与える。各ブロックへの射影像は対称二行行列の全空間であり、
その左右の同一視はスカラー以外一意である。実際、左右変換がこの空間を保てば全ての階数1平方の像も対称となり、
二つの変換は全ての射影直線上の点で一致する。従って各孤立点 $\operatorname{diag}(1,-1)$ は、合同の意味で固有の等方二直線を与える。

各等方対での評価を $U_\pm=A+B\pm2x$, $V_\pm=C+D\mp2x$ と取る。
$V_++V_-=-(a/b)(U_++U_-)$、$V_+-V_-=-(U_+-U_-)$ である。
このグラフ行列の対角積と反対角積の比は $((a+b)/(a-b))^2$ である。行列の行・列の個別倍率と置換、および
二平面交換による逆行列でこの比は逆数までしか変わらない。従って $r'^2=r^2$ または $r'^2=1/r^2$ が必要である。

後者を調べる。定理\ref{thm:transverse-doublet-isomorphism}の四次交代収縮を $D$、勾配を $G$ とする。
この規格化では $D=64(3r^4-5r^2+6)/(r^2-1)^2$。分子と逆順多項式は互いに素なので、$D=0$ は逆数と同型にならない。
$D\ne0$ のとき $L=W_{01\mid23}(G/D)_{01\mid23}^{t}$ を作る。
$k=2,3$ で $\operatorname{tr}(L^k)(r)-\operatorname{tr}(L^k)(1/r)$ の既約分子の最大公約数は定数倍を除いて $r^4-1$ である。
よって逆数と一致するのは最初から $r^2=1/r^2$ の場合だけである。二対の交換で $r$ の符号は変えられ、
$a+b=1/2$ から多重集合が決まる。
\end{proof}

\begin{proof}[異方向の二重色と一欠色曲線からの分離]
反対階数で三薄色の対称族を分離でき、捻りと異方向族の共通候補は階数16のみである。
その異方向表示の関係空間は、同じ行符号変更を施すと
\[
 \begin{pmatrix}
 ps+t&u&y&z\\u&qs+t&z&y\\y&-z&s+pt&v\\-z&y&v&s+qt
 \end{pmatrix}
\]
となる。$s=t=0$ の射影三空間内の二次曲面 $uv=y^2-z^2$ では全ての三次小行列式が消える。
三薄色の捻りの関係空間にはこの次元の低階数集合がなく、無印でも非同型である。

一欠色曲線が周期1となるのは $X=(t+1/t)^2=\pm4$ のみである。
反対切断の行列式は色0,1で $t^2(t^4+6t^2+1)/(t^2+1)$、色2,3で $-t^2(t^2-1)^2/(t^2+1)$。
$X=4$ は $t=1$ で代表され、左核の一般行列の行列式は基底を選べば
$(z_0^2-z_2^2)(z_1^2+z_3^2)/2$ となる。反対階数は12であり、この四相異平面は
異方向族の階数12の二境界の二重平面対および二つの既約二次錐と異なる。
$X=-4$ では消える二色の行列式が単純零で各階数は3、反対階数は14である。
共通候補は三薄色の捻りの $(-1,1,1,1)/2$ の置換だけである。両者とも $D\ne0$ で
$\operatorname{tr}(L^2)$ は順に $43/16,124/27$ となり分離される。
\path{computations/verify_three_lines_twisted_isomorphism.py} が余因子、例外直線の分子の最大公約数、
二次曲面、および欠色曲線の二つの交点候補を厳密計算する。
\end{proof}
